\section{Problem 154: counting divisible coefficients in Pascal's pyramid}
\label{sec:p154}

Problem 154 asks for one layer of Pascal's pyramid. The coefficient of
$x^a y^b z^c$ in $(x+y+z)^N$, where $a+b+c=N$, is
\[
  T(N;a,b,c)=\frac{N!}{a!b!c!}.
\]
Here $N=200000$. We must count the coefficients divisible by
$10^{12}=2^{12}5^{12}$. The implementation never constructs a coefficient
that large. It counts factors of 2 and 5 instead.

The organization of this explanation follows Tam Nguyen's
\href{https://tamnd.github.io/brain/practice/euler/154/}{Project Euler 154
note}: derive the valuation test, reduce the search by symmetry, and check the
logic on a smaller layer. The code and the two-stage optimization below come
from this repository.

\subsection{Replace divisibility with two valuation tests}

For a prime $p$, let $v_p(m)$ be the exponent of $p$ in $m$. Then
\[
 v_p(T)=v_p(N!)-v_p(a!)-v_p(b!)-v_p(c!).
\]
The coefficient qualifies exactly when $v_2(T)\geq12$ and
$v_5(T)\geq12$. Both programs build prefix tables for $v_2(n!)$ and
$v_5(n!)$. The recurrence is small:
\[
 v_p(n!)=v_p((n-1)!)+v_p(n).
\]
After this $O(N\log N)$ setup, either valuation of a binomial coefficient is
an $O(1)$ subtraction.

\subsection{Split one trinomial coefficient into two binomial coefficients}

Set $i=b+c$ and $j=b$. Then $a=N-i$ and $c=i-j$, so
\[
 T(N;N-i,j,i-j)=\binom{N}{i}\binom{i}{j}.
\]
Valuations turn the product into a sum:
\[
 v_p(T)=v_p\!\left(\binom{N}{i}\right)
       +v_p\!\left(\binom{i}{j}\right).
\]
This factorization creates a useful fast path. If $\binom{N}{i}$ already has
at least twelve factors of both 2 and 5, then all $i+1$ splits of $i$ into
$j$ and $i-j$ qualify. The program adds the whole row without inspecting its
members.

For the remaining rows, exchanging $b$ and $c$ does not change the
coefficient. The inner loop examines half of the row and gives a nonsymmetric
pair weight two. When $b=c$, the pair has weight one.

\subsection{C implementation}

The C version stores the two prefix tables in flat arrays. Its helper computes
the valuation of a binomial coefficient.

\begin{lstlisting}
static unsigned choose154(const unsigned *factorial, unsigned n, unsigned k) {
    return factorial[n] - factorial[k] - factorial[n - k];
}

for (unsigned i = 0; i <= LAYER; ++i) {
    unsigned outer2 = choose154(twos, LAYER, i);
    unsigned outer5 = choose154(fives, LAYER, i);
    if (outer2 >= 12 && outer5 >= 12) {
        answer += i + 1;
        continue;
    }
    for (unsigned j = 0; j <= (i + 1) / 2; ++j) {
        if (outer2 + choose154(twos, i, j) >= 12 &&
            outer5 + choose154(fives, i, j) >= 12) {
            ++answer;
            if (j < i / 2) ++answer;
        }
    }
}
\end{lstlisting}

The arrays contain at most $N+1$ values. The code uses unsigned subtraction
only where the factorial valuation is monotone, so the result cannot become
negative.

\subsection{Rust implementation}

Rust expresses the same algorithm with slices and vectors. The loop bounds,
the row fast path, and the symmetry weights match the C version.

\begin{lstlisting}
fn choose(v: &[u32], n: usize, k: usize) -> u32 {
    v[n] - v[k] - v[n - k]
}

for i in 0..=N {
    let (a, b) = (choose(&twos, N, i), choose(&fives, N, i));
    if a >= 12 && b >= 12 {
        answer += (i + 1) as u64;
        continue;
    }
    for j in 0..=(i + 1) / 2 {
        if a + choose(&twos, i, j) >= 12
            && b + choose(&fives, i, j) >= 12 {
            answer += 1;
            if j < i / 2 { answer += 1 }
        }
    }
}
\end{lstlisting}

The Rust types state two bounds that C leaves implicit. Table indices use
\texttt{usize}, and the final count uses \texttt{u64}. The valuation entries fit
in \texttt{u32} by a wide margin.

\subsection{Cross-language comparison and missing ports}

C and Rust use the same mathematical decomposition and asymptotic work. C
manages the table allocation and release. Rust owns the tables through
\texttt{Vec<u32>} and releases them when \texttt{solve} returns. Neither version
uses the shared prime or combinatoric API because the two valuation tables are
specific to this problem.

The repository has no validated Gforth or Factor implementation for problem
154. Their scaffold files stay outside answer coverage. This report does not
replace those missing programs with pseudocode or reference-answer constants.

\subsection{Small checks catch symmetry errors}

For a reduced layer, enumerate every ordered triple $(a,b,c)$ with
$a+b+c=n$. Compare the direct multinomial valuation with the two-binomial
decomposition. Then compare the full ordered count with the weighted
half-row count. These checks detect the two common mistakes: dropping the
middle pair when $i$ is odd and counting the $b=c$ case twice.

The generated figures in Section~\ref{sec:visualatlas} show three views of the
same object. The coefficient surface explains why central entries grow fastest.
The parity plot shows what remains when divisibility is reduced to one factor
of 2. Problem 154 applies the stronger two-prime test at layer 200000.
